%% RePLab — Overleaf scientific short paper (CHAINS-style structure)
%% Upload this folder to Overleaf as a new project, or paste main.tex as the sole source.
%% Edit \author fields before submission.
\documentclass[11pt,a4paper]{article}

\usepackage[margin=1in]{geometry}
\usepackage[T1]{fontenc}
\usepackage[utf8]{inputenc}
\usepackage{lmodern}
\usepackage{microtype}
\usepackage{booktabs}
\usepackage{graphicx}
\usepackage{enumitem}
\usepackage{caption}
\usepackage{subcaption}
\usepackage{amsmath}
\usepackage{amssymb}
\usepackage{xcolor}
\usepackage{hyperref}
\usepackage{titlesec}
\usepackage{fancyhdr}
\usepackage{setspace}

\hypersetup{
  colorlinks=true,
  linkcolor=black,
  citecolor=blue!50!black,
  urlcolor=blue!60!black
}

\pagestyle{fancy}
\fancyhf{}
\fancyhead[L]{\small RePLab: HITL Multi-Agent Paper Reproduction}
\fancyhead[R]{\small Final Project Report}
\fancyfoot[C]{\thepage}
\renewcommand{\headrulewidth}{0.4pt}

\title{\textbf{RePLab:} A Human-in-the-Loop Multi-Agent System\\
for Finding, Assessing, and Reproducing Papers with Public Code}
\author{
  %% --- edit before submission ---
  Oscar\\
  Department of Computer Science\\
  %% University / Program
  Master's Web Development / Agentic AI Final Project\\
  \texttt{[email@university.edu]}
}
\date{}

\begin{document}
\maketitle

\begin{abstract}
Large Language Model (LLM) agents can accelerate literature search and experiment
prototyping, yet end-to-end ``reproduce the paper'' prompts often yield
non-executable text, unsupported claims, and unsafe or wasteful host actions.
This paper presents \textbf{RePLab} (Paper-with-Code Reproduction Lab), a
Human-in-the-Loop (HITL) multi-agent environment that restricts candidates to
research papers with \emph{public code repositories} and decomposes reproduction
into inspectable stages: paper retrieval with code validation, user selection,
deterministic-plus-LLM feasibility analysis, permission-gated sandboxed
execution, and parameter tweaking with comparison. Each stage emits typed
artifacts (candidates, feasibility reports, run logs, metrics) under an explicit
verdict policy (\texttt{IMPOSSIBLE}, \texttt{RISKY\_BUT\_POSSIBLE}, \texttt{READY}).
A hybrid deployment model exposes search and analysis on serverless hosts while
keeping clone/install/run local under Docker. We situate RePLab relative to
prior agentic research systems, specify an evaluation protocol under matched
budgets, and report a design walkthrough of the implemented LangGraph/Gradio
system. RePLab treats agents as coordinated, governable components---not opaque
replacements---for accountable scientific code reproduction.
\end{abstract}

\noindent\textbf{Keywords:}
Agentic AI; Multi-Agent Systems; Human-in-the-Loop; Scientific Reproducibility;
Papers with Code; LangGraph; Docker Sandboxing; Digital Society Systems.

%% ========================================================================
\section{Introduction and Motivation}
\label{sec:intro}
%% ========================================================================

Scientific workflows increasingly rely on LLM assistants for literature review,
summarization, and coding support. Prompt-only pipelines, however, frequently
provide limited provenance, weak reproducibility, and few safeguards against
unsupported claims or high-risk actions on the researcher's machine~\cite{chains2026,react}.
In particular, asking a single model to ``reproduce this paper'' tends to
produce \emph{paper-shaped text}: fluent narratives that look like experimental
results but are not grounded in an executable public codebase.

Reproducibility in machine learning and related fields is already fragile when
code, data, and environment pins are incomplete~\cite{pineau2021}. Agentic
systems amplify both opportunity and risk: they can search, clone, install, and
run software autonomously, which makes \emph{governance}---when the human must
approve resource-consuming or irreversible steps---a first-class design
requirement rather than an afterthought~\cite{autogen,chains2026}.

This work introduces \textbf{RePLab}, a gated multi-agent lab with three design
commitments:
\begin{enumerate}[leftmargin=*,itemsep=2pt]
  \item \textbf{Code-anchored candidates.} Only papers with validated public
        GitHub (or equivalent) repositories enter the working set.
  \item \textbf{Honest feasibility before execution.} A structured verdict
        stops impossible jobs and warns on soft risks before any clone/run.
  \item \textbf{Explicit HITL gates.} The human selects the paper and must
        approve sandboxed reproduction; agents do not silently consume disk,
        CPU/GPU, or network quotas.
\end{enumerate}

RePLab is implemented as a Gradio application with FastAPI packaging for hybrid
cloud deployment, LangGraph-based research scouts, and a staged reproduction
crew (\texttt{PaperFinder}, \texttt{FeasibilityAnalyst},
\texttt{ExperimentRunner}, \texttt{ParamTweaker}). The system is intended for
student and laboratory settings where transparency and safety matter as much as
automation.

%% ========================================================================
\section{Background and Related Work}
\label{sec:background}
%% ========================================================================

\subsection{Definitions}
\begin{description}[style=nextline,leftmargin=1.2cm,itemsep=2pt]
  \item[Agentic AI]
    Specialized LLM agents that plan, call tools, and iterate until a goal is
    met~\cite{react}.
  \item[Artifact-driven workflow]
    Intermediate outputs (plans, notes, reports, logs) are persisted as typed,
    inspectable artifacts rather than discarded chat turns~\cite{metagpt,chains2026}.
  \item[Human-in-the-Loop (HITL)]
    Humans retain authority to select, approve, decline, or stop agent actions
    at critical checkpoints~\cite{autogen,chains2026}.
  \item[Deterministic verification]
    Rule-based (non-generative) checks---e.g., entrypoint file existence,
    dependency presence---that constrain or override LLM suggestions.
\end{description}

\subsection{Multi-agent collaboration patterns}
Prior systems illustrate recurring organizational patterns that inform RePLab:
\begin{itemize}[leftmargin=*,itemsep=2pt]
  \item \textbf{Orchestrator + specialists} (Magentic-One): a lead agent plans
        and reassigns browser/file/code tools~\cite{magenticone}.
  \item \textbf{Role chat / SOPs} (ChatDev, MetaGPT): software-company roles
        and document handoffs~\cite{chatdev,metagpt}.
  \item \textbf{Conversational teams + HITL} (AutoGen): multi-agent meetings
        with optional human proxies~\cite{autogen}.
  \item \textbf{Task crews} (CrewAI): explicit role--task pipelines~\cite{crewai}.
  \item \textbf{Research swarms with manifests} (CHAINS): planning, retrieval,
        gap finding, micro-experiments, drafting, critique, and deterministic
        verification under a run manifest~\cite{chains2026}.
\end{itemize}
RePLab adopts CHAINS' emphasis on HITL, artifacts, and non-generative checks,
but specializes the product goal: \emph{code-backed reproduction} with Docker
sandboxing, rather than end-to-end paper drafting from knowledge gaps.

\subsection{Papers with code}
Linking publications to official implementations improves the chance that
reproduction is even possible~\cite{paperswithcode}. RePLab therefore treats
public repositories as a hard eligibility filter, not a nice-to-have metadata
field.

%% ========================================================================
\section{Research Questions}
\label{sec:rqs}
%% ========================================================================

\begin{description}[style=nextline,leftmargin=1.35cm,itemsep=4pt]
  \item[RQ1 (Operationalization)]
    How effectively can a staged agent pipeline turn a topic query into a
    \emph{runnable} paper-with-code candidate set with explicit feasibility
    labels?
  \item[RQ2 (Transparency)]
    Does an artifact-driven, verdict-gated workflow improve inspectability and
    failure localization compared with a single LLM prompt that claims to
    reproduce a paper?
  \item[RQ3 (Trust and safety)]
    How do HITL checkpoints (paper selection; permission to clone/install/run)
    influence prevention of impossible or unsafe executions and user control
    over resource use?
\end{description}

%% ========================================================================
\section{System Architecture}
\label{sec:arch}
%% ========================================================================

Figure~\ref{fig:arch} summarizes the deployment and control planes.

\begin{figure}[t]
\centering
\fbox{\parbox{0.92\linewidth}{\centering\small
\textbf{UI (Gradio)} --- secure API-key field; Reproduction Lab; Research Scout\\[4pt]
$\downarrow$\\[2pt]
\textbf{ASGI (FastAPI)} --- hybrid Vercel packaging; security headers\\[4pt]
$\downarrow$\\[2pt]
\begin{tabular}{@{}c@{\hspace{1.2em}}c@{}}
\textbf{RePLab crew} & \textbf{Scout graphs} \\
Finder $\rightarrow$ Analyst & Single ReAct \\
$\rightarrow$ Runner $\rightarrow$ Tweaker & Planner--Researcher--Writer--Critic \\
\end{tabular}\\[4pt]
$\downarrow$\\[2pt]
\textbf{External services:} OpenAI; OpenAlex; arXiv; Semantic Scholar; GitHub\\[2pt]
\textbf{Local sandbox:} Docker (\texttt{python:3.11-slim}); \texttt{runs/} artifacts
}}
\caption{RePLab system architecture: UI orchestration, agent crews, external APIs, and local Docker sandbox.}
\label{fig:arch}
\end{figure}

\subsection{Technology stack}
\begin{itemize}[leftmargin=*,itemsep=2pt]
  \item \textbf{Frontend / UX:} Gradio with live ``Reasoning:'' status streaming
        (no opaque progress-only bars).
  \item \textbf{Packaging:} FastAPI mounts the Gradio app for ASGI deploy;
        \texttt{vercel.json} configures duration/memory and basic security
        headers.
  \item \textbf{Orchestration:} staged Python agents in \texttt{replab/};
        LangGraph for optional research scouts.
  \item \textbf{Schemas:} Pydantic models for candidates, feasibility,
        entrypoint plans, run results, and tweak plans.
  \item \textbf{Secrets:} request-scoped API-key override (password field);
        keys are not written to \texttt{.env} from the UI, run folders, or
        reports; status text is redacted against key-shaped substrings.
\end{itemize}

\subsection{Hybrid deploy matrix}
Serverless environments can host search and feasibility, but not long-running
Docker builds. RePLab therefore uses an explicit hybrid policy
(Table~\ref{tab:hybrid}).

\begin{table}[ht]
\centering
\caption{Hybrid capability matrix.}
\label{tab:hybrid}
\begin{tabular}{@{}lll@{}}
\toprule
\textbf{Stage} & \textbf{Local} & \textbf{Serverless (e.g., Vercel)} \\
\midrule
Find papers + feasibility & Yes & Yes \\
HITL select / permission UI & Yes & Yes \\
Docker clone / install / run & Yes & No (emit local commands) \\
Parameter re-runs & Yes & No \\
\bottomrule
\end{tabular}
\end{table}

%% ========================================================================
\section{Multi-Agent Workflow}
\label{sec:workflow}
%% ========================================================================

\subsection{Reproduction Lab crew}
The primary product path is a four-agent (plus tooling) pipeline:

\begin{description}[style=nextline,leftmargin=1.35cm,itemsep=3pt]
  \item[PaperFinder]
    Queries \textbf{OpenAlex} as primary source (arXiv-linked works with GitHub
    signals), then arXiv / Semantic Scholar when needed; validates repositories
    via the GitHub API; screens feasibility during search; ranks
    \texttt{READY} before \texttt{RISKY}, then by year and stars.
  \item[FeasibilityAnalyst]
    Emits a structured \texttt{FeasibilityReport}: summary, Python/CUDA/package
    needs, size/runtime estimates, alarms, entrypoint, and verdict. Deterministic
    repo signals (README, requirements, tree, CUDA hints) are merged with optional
    LLM structuring; hallucinated entrypoints that do not exist on disk (via API
    tree) are rejected.
  \item[ExperimentRunner]
    After permission: stages under \texttt{runs/in\_progress/}, shallow-clones,
    prepares an entrypoint command (AST argparse/click parsing, README demos,
    Streamlit headless handling), runs inside Docker with memory/CPU limits and
    timeouts, retries once with deterministic then LLM repair suggestions, and
    archives to \texttt{runs/successful/} or \texttt{runs/failed/}.
  \item[ParamTweaker]
    Proposes a small editable parameter set, re-invokes the runner with
    overrides, and formats a baseline-vs-tweaked comparison table.
\end{description}

\textbf{HITL gate~1 --- selection.} The user chooses a paper from a screened
table (title, year, verdict, entrypoint, repo).\\
\textbf{HITL gate~2 --- permission.} Clone/install/run proceeds only if the
user checks an explicit approval control and the verdict is not
\texttt{IMPOSSIBLE}.

\subsection{Verdict policy}
\begin{itemize}[leftmargin=*,itemsep=2pt]
  \item \texttt{IMPOSSIBLE}: missing runnable entrypoint, stub-only repo,
        private required data, or no reduced demo path for extreme compute ---
        \emph{stop}.
  \item \texttt{RISKY\_BUT\_POSSIBLE}: soft limits (CUDA likely, large downloads,
        outdated pins, long runtime) --- \emph{warn and require permission}.
  \item \texttt{READY}: plausible minimal path --- \emph{confirm and run}.
\end{itemize}

\subsection{Entrypoint intelligence}
A dedicated preflight module sanitizes commands, resolves script paths, parses
CLI arguments from source, fills required positionals from documentation
examples, and avoids trusting raw LLM command strings. This deterministic layer
is RePLab's analogue of CHAINS' non-generative verifier, applied to
\emph{execution readiness} rather than citation structure.

\subsection{Optional Research Scout}
For literature synthesis (not reproduction), RePLab retains:
\begin{itemize}[leftmargin=*,itemsep=2pt]
  \item a \textbf{single ReAct} scout (\texttt{web\_search}, \texttt{fetch\_page});
  \item a \textbf{team} graph: Planner $\rightarrow$ Researcher $\leftrightarrow$
        tools $\rightarrow$ Writer $\rightarrow$ Critic (retry/ok) $\rightarrow$
        finalize.
\end{itemize}
These paths share the secure API-key plumbing and produce Markdown reports under
\texttt{outputs/}.

\subsection{End-to-end control flow}
\begin{verbatim}
START
  -> PaperFinder (OpenAlex-first; GitHub validated; feasibility screen)
  -> USER: select paper                         # HITL #1
  -> FeasibilityAnalyst (typed report + verdict)
       |-- IMPOSSIBLE -> explain -> END
       +-- RISKY / READY -> USER: approve run?  # HITL #2
              |-- no  -> END
              +-- yes -> ExperimentRunner (Docker, local)
                          -> ParamTweaker (tweak -> re-run -> compare)
                          -> archive artifacts -> END
\end{verbatim}

%% ========================================================================
\section{Evaluation Protocol and Quality Metrics}
\label{sec:metrics}
%% ========================================================================

Following CHAINS' practice of defining measurable proxies before claiming
superiority~\cite{chains2026}, we specify a protocol comparing RePLab to an
\textbf{LLM-only baseline} (single prompt: search + summarize + ``reproduce'')
under matched model, token, and wall-clock budgets.

\subsection{Proposed metrics}
\begin{description}[style=nextline,leftmargin=1.4cm,itemsep=3pt]
  \item[WCR (Workflow Completion Rate)]
    Fraction of topics for which the pipeline reaches a terminal state with all
    required artifacts present (candidate list, feasibility report; and, if
    permitted, run archive).
  \item[ABI (Artifact Bundle Integrity)]
    Fraction of required typed fields present and schema-valid (title, repo URL,
    verdict, entrypoint or explicit impossibility reason, run log path when
    executed).
  \item[FLA (Feasibility Label Agreement)]
    Agreement between system verdicts and a human rater on a held-out set
    (Cohen's $\kappa$), with disagreement analysis on entrypoint and data
    availability.
  \item[RCP (Run Completion under Permission)]
    Among user-approved, non-\texttt{IMPOSSIBLE} cases, fraction that exit with
    success criteria (process exit code 0 or documented soft-success, e.g.,
    Streamlit smoke boot).
  \item[HITL-SE (HITL Safety Efficiency)]
    Rate of prevented high-risk actions: attempts blocked by
    \texttt{IMPOSSIBLE} or missing permission, plus false-allow incidents
    (should have been blocked).
  \item[GAS-lite (Grounding Accuracy Score, lite)]
    Sampled audit: do feasibility claims (entrypoint, CUDA need, key packages)
    match repository evidence?
\end{description}

\subsection{Procedure}
For $N$ topics spanning NLP, vision, and time-series demos:
(1)~run baseline and RePLab under matched budgets;
(2)~score WCR/ABI automatically from artifacts;
(3)~double-annotate FLA/GAS-lite on a subsample;
(4)~record latency, token use, and HITL decision times.
This protocol is specified for reproducibility of the \emph{evaluation design};
large-scale numerical benchmarking remains future coursework/lab work
(Section~\ref{sec:limits}).

%% ========================================================================
\section{Demonstration and Qualitative Findings}
\label{sec:results}
%% ========================================================================

We demonstrate RePLab through an implemented end-to-end system and smoke checks
rather than claiming a finished multi-topic benchmark.

\subsection{Implemented capabilities}
\begin{itemize}[leftmargin=*,itemsep=2pt]
  \item OpenAlex-first retrieval with rate-limit aware backoff and optional
        Semantic Scholar / GitHub tokens.
  \item Search-time feasibility screening so the UI preferentially shows
        runnable candidates.
  \item Live reasoning traces in the UI for finder, analyst, runner, and
        tweaker stages.
  \item Docker-bounded execution with archival under \texttt{runs/successful}
        and \texttt{runs/failed}.
  \item Success narration that explains, in plain language, what ran and why
        arguments matter.
  \item Integration smoke script validating finder/analyst paths and a minimal
        Docker invocation when the engine is available.
\end{itemize}

\subsection{Qualitative observations (design walkthrough)}
\begin{enumerate}[leftmargin=*,itemsep=3pt]
  \item \textbf{Paper-shaped text vs.\ executable path.}
        LLM-only baselines readily invent train commands; RePLab's file-existence
        checks and AST CLI parsing reduce command hallucination before Docker.
  \item \textbf{Failure localization.}
        Separating find / assess / run / tweak makes it obvious whether a
        session failed due to retrieval, feasibility, missing permission, or
        sandbox errors---unlike a single chat transcript.
  \item \textbf{Governance value of HITL.}
        Permission gating prevents accidental multi-gigabyte clones and long
        installs when the user only intended to browse candidates.
  \item \textbf{Hybrid realism.}
        Cloud UIs are useful for discovery; honest ``run locally'' messaging
        avoids pretending serverless hosts can replace laboratory machines.
\end{enumerate}

These observations align with CHAINS' findings that structured, artifact-centric
workflows improve auditability relative to black-box generation~\cite{chains2026},
while specializing the claim to \emph{code reproduction}.

%% ========================================================================
\section{Limitations and Future Work}
\label{sec:limits}
%% ========================================================================

\begin{itemize}[leftmargin=*,itemsep=2pt]
  \item \textbf{Preliminary evaluation.} Metrics in Section~\ref{sec:metrics}
        are specified; full matched-budget scoring across diverse domains is
        not yet reported.
  \item \textbf{Semantic truthfulness.} Deterministic checks and GAS-lite do
        not guarantee that a successful smoke run equals a faithful paper
        result (metrics extraction is heuristic).
  \item \textbf{Environment coverage.} The default image targets CPU Python
        demos; CUDA-heavy repositories are often labeled \texttt{RISKY} and may
        still fail without host GPU passthrough and specialized images.
  \item \textbf{Data and licenses.} Private datasets, gated downloads, and
        unclear licenses remain hard stops or manual steps.
  \item \textbf{Scale and comparison.} Broader benchmarking against Magentic-One,
        AutoGen, CrewAI, and CHAINS-style research swarms is future work.
\end{itemize}

\noindent\textbf{Future directions:} stronger semantic grounding of feasibility
claims; richer dataset tooling; GPU-aware runners; signed run manifests with
content hashes (closer to CHAINS' telemetry); and user studies on warning
usefulness and trust.

%% ========================================================================
\section{Conclusion}
\label{sec:conclusion}
%% ========================================================================

RePLab frames AI agents not as opaque end-to-end replacements for scientific
reproduction, but as \textbf{coordinated components in a human-governed,
code-anchored workflow}. By combining specialized agents, typed artifacts,
deterministic entrypoint verification, Docker sandboxing, and explicit HITL
gates, the system prioritizes transparency, traceability, and controlled
execution---prerequisites for responsible AI assistance in laboratory and
classroom settings. The core lesson mirrors CHAINS while narrowing the task:
\emph{accountable reproduction begins with public code, honest feasibility, and
permission---not with a single generative claim that an experiment ``worked.''}

\section*{Acknowledgments}
This report was prepared as a Master's final project artifact for an agentic AI /
web development coursework setting. The design is informed by CHAINS and related
multi-agent research systems cited herein.

%% ========================================================================
\begin{thebibliography}{99}

\bibitem{chains2026}
W.~Villalobos, Y.~Kumar, J.~Marchena, J.~J.~Li, and D.~Kruger.
``Collaborative Human-`AI ageNt Swarm' (CHAINS) for Conducting Scientific Research.''
In \emph{ICDS 2026}, ThinkMind / IARIA, Venice, Italy, 2026, pp.~8--15.
\url{https://www.thinkmind.org/library/ICDS/ICDS_2026/icds_2026_1_30_10034.html}

\bibitem{react}
S.~Yao, J.~Zhao, D.~Yu, N.~Du, I.~Shafran, K.~Narasimhan, and Y.~Cao.
``ReAct: Synergizing Reasoning and Acting in Language Models.''
In \emph{ICLR}, 2023.
\url{https://arxiv.org/abs/2210.03629}

\bibitem{pineau2021}
J.~Pineau et~al.
``Improving Reproducibility in Machine Learning Research (A Report from the NeurIPS 2019 Reproducibility Program).''
\emph{Journal of Machine Learning Research}, 22(164):1--20, 2021.

\bibitem{magenticone}
A.~Fourney et~al.
``Magentic-One: A Generalist Multi-Agent System for Solving Complex Tasks.''
Microsoft Research, 2024.
\url{https://aka.ms/magentic-one}

\bibitem{chatdev}
C.~Qian et~al.
``ChatDev: Communicative Agents for Software Development.''
arXiv:2307.07924, 2023.

\bibitem{metagpt}
S.~Hong et~al.
``MetaGPT: Meta Programming for A Multi-Agent Collaborative Framework.''
arXiv:2308.00352, 2023.

\bibitem{autogen}
Q.~Wu et~al.
``AutoGen: Enabling Next-Gen LLM Applications via Multi-Agent Conversation.''
arXiv:2308.08155, 2023.

\bibitem{crewai}
CrewAI documentation.
\url{https://docs.crewai.com/}

\bibitem{paperswithcode}
Papers with Code.
\url{https://paperswithcode.com/}

\end{thebibliography}

\end{document}
